\chapter{工业数据工程与实验设计}

\begin{quote}\small
\textit{本章摘要。}本章处理模型之前和模型之外的必要问题：数据从哪里来、时间如何对齐、缺失与异常如何解释、实验如何复现以及系统如何验收。本章结合数据生命周期讨论实验设计，是因为不可追溯的数据会使性能提升无法归因，也无法在部署后修复。
\end{quote}

前两章要求独立评价、清楚的目标分布和不泄漏的模型选择，但这些条件不会由一条训练命令自动满足。它们首先是数据工程问题：一条温度记录在预测时刻是否已经到达，设备停机究竟是故障还是计划维护，同一运行过程是否被重复算成独立证据。本章把统计假设翻译成可以执行和追溯的数据规则。

表格里的一行数据，往往经过截取窗口、补齐缺失值和匹配工单才得到。每一步都可能改变任务：用停机后的读数补停机前的缺失，就让模型提前看到了未来；把计划维护当成故障，则改错了标签。因此要先确定预测时刻和目标，再决定怎样清洗。最后按记录当时最终可用的时刻回放，并通过不直接控制设备的影子运行，检查线上是否能得到同样的输入和结果。

本章从一条记录、一个预测时刻和一个标签开始，不要求先掌握数据库或流式框架。阅读时先画出输入截止、标签观察截止和训练截止三条时间线，再检查每个字段何时可见；频谱与删失似然分别服务于信号特征和未完成寿命记录，可在相应任务出现时细读。

\section{从原始记录到可学习样本}

% auxiliary-resource-guide
本章末的“辅助学习材料”列出与核心方法对应的讲解和实践入口，可在阅读相关推导后，按所列小节对照学习。

\begin{figure}[htbp]
\centering
\begin{tikzpicture}[node distance=0.7cm and 0.65cm]
 \node[idi input, minimum width=2.0cm] (event) {事件发生\\$t_{event}$};
 \node[idi process, right=of event, minimum width=2.0cm] (collect) {采集与传输\\$t_{collect}$};
 \node[idi state, right=of collect, minimum width=2.0cm] (usable) {系统可用\\$t_{avail}$};
 \node[idi output, right=of usable, minimum width=2.0cm] (predict) {预测时刻\\$t$};
 \draw[idi arrow] (event)--(collect); \draw[idi arrow] (collect)--(usable); \draw[idi arrow] (usable)--(predict);
 \draw[idi feedback] (predict.south) .. controls +(0,-0.42) and +(0,-0.42) .. node[below,font=\scriptsize]{只允许使用 $a_j\leq t$ 的可用记录} (usable.south);
\end{tikzpicture}
\caption{工业样本的时间边界至少包括事件、采集和可用时刻；预测输入不能回看未来才进入系统的信息。}
\label{fig:data-availability-timeline}
\end{figure}
图\ref{fig:data-availability-timeline}将事件发生、采集和系统可用时刻分开，强调构造预测输入时应检查记录何时可查询，而不只是何时产生。
\paragraph{数据生命周期。}
工业数据从传感器、控制器、日志系统、工单和实验室检测中产生。一个可学习的数据集必须明确采集时间、设备标识、工况、单位、采样频率、缺失含义和标签来源。数据字典应把“缺失”“零值”“未采集”和“传感器故障”区分开，否则模型会把数据质量问题当作物理规律。

\paragraph{时间对齐与窗口化。}
不同传感器的时间戳可能存在时钟偏移、延迟和不同采样频率。对齐时需要选择参考时钟，记录插值方法，并保留原始时间间隔。长度为$L$、步长为$s$的滑动窗口为
\begin{equation}
 X_t=[x_{t-L+1},\ldots,x_t],
 \qquad y_t=\text{在未来区间定义的标签}.
\end{equation}
标签必须明确预测时刻和观察截止时刻，避免把未来维修记录泄漏到当前样本。

\paragraph{用信息集合推导合法样本。}
设第$j$条记录为$(v_j,e_j,a_j)$，其中$v_j\in\R^d$、$e_j$为事件时刻、$a_j$为该版本首次可用时刻。预测时刻$t$的信息集合记作$\mathcal F_t=\sigma\{(v_j,e_j,a_j):a_j\leq t\}$；合法特征必须由$\mathcal F_t$中的记录计算。若使用最近$W$小时的观测，输入为
\begin{equation}
 x_t=\phi\{v_j:t-W<e_j\leq t,\ a_j\leq t\}.
\end{equation}
两个条件缺一不可：事件发生得早，不代表当时已被系统接收。回填的维修结论和修正后的传感器值必须保留版本时间，否则历史回放会读到当时不可见的信息。

定义$t$之后下一次目标故障时刻为$T_f$，没有后续目标故障时取$T_f=\infty$，则$y_t=\mathbb 1\{t<T_f\leq t+H\}$。阴性标签必须等到$t+H$且完成报告延迟后才能确认；若运行记录提前结束，应标为未确定而非零。在线输入的合法性由$t$判断，训练标签是否成熟则由训练截止时间$c$判断。这是两个不同的时间检查。

对等间隔序列$x_r\in\R^d$，窗口矩阵可约定$X_t\in\R^{L\times d}$，第$r$行为$x_{t-L+r}^{\mathsf T}$。不规则时间序列不能仅用行数代表持续时间，需要同时保留$\Delta t$或按真实时间区间构造窗口。在线前向填补只使用过去最近值；线性插值若用到了$t$以后的点就违反$\mathcal F_t$可测性。

\paragraph{用一条迟到记录检查两道时间门。}
$\sigma\{\cdots\}$表示由已有记录生成的信息集合，不是标准差；“可测”在此可先读为“只靠这些已知记录就能算出”。例如预测时刻为10时，窗口宽度$W=2$小时：9时50分产生、10时05分才入库的读数不能用于10时预测；9时40分产生、9时45分入库的读数可以。若目标是未来$H=1$小时故障，10时30分故障属于正例，但该结论只能在故障确认后作为训练标签，不能放进10时的输入。假如记录到10时40分就结束且尚无故障，则无法确认11时之前不会故障，应保留未确定标记。

窗口长度$L$计数的是采样点，$W$计数的是实际时间，二者不要混用。例如每分钟采样两个通道，$L=3$的窗口是$3\times2$矩阵；步长$s=1$意味着下一窗口只前移一个采样点，两个窗口共享两行，而不是得到两次独立运行。

\paragraph{划分间隔由哪些窗口决定。}
若训练标签覆盖$(t,t+H]$，而验证从$v$开始，为避免训练标签读取验证时期，至少要求训练样本$t+H\leq v$；这不等于输入完全不重叠。若还要求两组原始输入窗口无交叠，则训练最后时刻与验证首时刻之间要留足$W$的间隔。业务上允许共享已知历史时，可以不要求后一条件，但应明确估计的是同设备未来能力，而非独立设备泛化。间隔长度由信息集合决定，不能机械地设置为某个固定百分比。

\paragraph{缺失值与异常值。}
缺失机制可粗略分为完全随机缺失、条件随机缺失和非随机缺失。均值填补简单但会压低方差；前向填补适合连续状态但可能延续故障；模型填补可以利用多变量关系却会引入额外不确定性。建议同时加入缺失指示变量，并对不同缺失机制进行敏感性分析。

异常值不应一律删除。尖峰可能是真实冲击，也可能是通信错误。应区分物理约束异常、统计异常和任务异常，并将修正规则记录为可复现的数据版本。

\paragraph{填补究竟保留了什么。}
令$M=1$表示变量$V$被观测到。完全随机缺失要求$M$与完整数据独立；条件随机缺失允许$M$依赖已观测变量，但给定它们后不再依赖缺失值本身；非随机缺失不满足后者。仅从已有表格一般不能检验最后两种机制的区别，因此需要采集知识或敏感性假设。

设$M\sim\operatorname{Bernoulli}(r)$与$V$独立，$\E V=\mu$、$\Var(V)=\sigma^2$，用真实均值填补得到$\tilde V=MV+(1-M)\mu$。则
\begin{equation}
 \E\tilde V=\mu,\qquad
 \Var(\tilde V)=\E[M(V-\mu)^2]=r\sigma^2.
\end{equation}
即使均值完全已知，填补也会把方差压缩，故把填补值当作真实测量会低估不确定性。若用其他已观测变量$Z$，平方损失下的最优点填补为$g(Z)=\E[V\mid Z]$，因为
\begin{equation}
 \E[(V-g(Z))^2\mid Z]
 =\Var(V\mid Z)+(\E[V\mid Z]-g(Z))^2.
\end{equation}
点填补只能保留条件均值，不能恢复条件方差。加入缺失掩码和“距上次真实观测的时间”使模型能区分测量与代入，但仍不能识别未观测的非随机缺失机制。所有填补器及均值、尺度都只在训练折拟合，验证和部署仅应用这些参数。

\paragraph{缺失不是零的一个手算例子。}
三条温度记录为$20,\text{缺失},24$，训练侧均值为22时，填补结果为$(20,22,24)$，同时保存观测掩码$M=(1,0,1)$，与上文$M=1$表示“已观测”的约定一致。若实现需要“缺失为一”的指示，应明确另定义为$1-M=(0,1,0)$，其中每项对应一条记录，避免把两种掩码混用。若独立验证记录为$(100,\text{缺失})$，仍应代入训练均值22，而不是从验证记录重新算均值100。均值填补并未断言缺失时真实温度为22，只是为算法提供可计算输入。

\paragraph{特征工程与表示学习。}
旋转机械常用时域统计量、频域能量和包络谱特征；过程数据常用滞后量、变化率、累计量和工况归一化值。特征工程把领域知识写入输入，深度表示学习则从原始信号寻找有用变换。两者可以并行比较，不应预先假定端到端网络一定优于物理启发特征。

\paragraph{从采样信号推导可比特征。}
对等间隔实信号$v_0,\ldots,v_{L-1}$，采样间隔$\Delta$、采样率$f_s=1/\Delta$，均值和均方根分别为$\bar v=L^{-1}\sum_rv_r$、$v_{\rm rms}=\sqrt{L^{-1}\sum_rv_r^2}$。两者满足$v_{\rm rms}^2=\bar v^2+L^{-1}\sum_r(v_r-\bar v)^2$，所以去均值前后的Root Mean Square (RMS)描述不同物理量。

取离散傅里叶变换$V_k=\sum_{r=0}^{L-1}v_r\exp(-2\pi\mathrm{i}kr/L)$。离散复指数的正交性给出
\begin{equation}
 \sum_{k=0}^{L-1}|V_k|^2
 =\sum_{r,s}v_rv_s\sum_k e^{-2\pi\mathrm{i}k(r-s)/L}
 =L\sum_{r=0}^{L-1}v_r^2.
\end{equation}
因此$L^{-2}\sum_k|V_k|^2$与时域均方一致。频带能量可以在目标频率索引上求和，但必须说明是双边频谱还是单边频谱；单边谱除直流和偶数长度时的奈奎斯特点外要计入共轭频率的能量。频率间隔为$f_s/L$，不是增加零填充后就获得更多物理分辨率。

该计算以等间隔采样为前提，降采样前需低通抑制混叠，有限窗口与加窗会改变谱泄漏和能量归一化。转速变化时固定频带也可能不再跟随故障阶次。这说明特征不是纯数学的列变换，其单位、采样机制和工况含义都属于实验设置。

\paragraph{先用四个点理解频域单位。}
若$f_s=4$赫兹、$L=4$，双边离散频率间隔为1赫兹。信号$(1,-1,1,-1)$的离散变换只有$V_2=4$非零，对应2赫兹的奈奎斯特频率；时域均方为1，频域计算为$4^2/4^2=1$。单边谱中这一端点不能再翻倍，否则能量错成2。对恒定信号$(2,2,2,2)$，均值和RMS都是2，但去均值后的RMS是零，分别表达静态水平与围绕均值的波动。这里$\mathrm{i}^2=-1$是虚数单位，$|V_k|^2$是复数幅度的平方，不是直接将复数平方。

\section{从数据处理到可信实验}
数据依赖、反馈回路和配置变化会在长期运行中积累系统风险，机器学习系统的技术债务研究对此给出了系统性讨论\cite{sculley2015debt}。以下版本与来源记录要求用于追溯预测依据，不能只记录模型权重。
\paragraph{实验矩阵。}
高质量实验应一次只改变一个主要因素。建议建立实验矩阵，至少包含数据切分、输入窗口、模型、损失、超参数、随机种子、硬件、训练时间、参数量、延迟和评价指标。主结果应包含简单基线、当前方法、消融实验、跨域测试和失败案例。

\paragraph{可复现性。}
代码版本、数据版本、环境文件、配置文件和随机种子共同定义一次实验。随机种子不能保证不同硬件完全一致，但可以帮助发现结果是否依赖偶然初始化。日志应保存训练曲线、最佳检查点、验证预测、错误样本和配置哈希。

\paragraph{能复现一次分数，不等于能解释改善。}
假设新模型误差由5降为4，但同时换了设备切分、删掉困难样本并调长窗口，就不能把下降归因于网络结构。应先固定设备清单、指标单位与预处理，再比较原模型和新模型；若必须改标签规则，应把它作为另一实验因素并重新评价双方。“检查点”是某一训练时刻的参数文件，“配置哈希”是配置内容的指纹，帮助判断两次运行是否使用同一设置，但都不能代替保存原始配置和数据来源。

\section{从离线证据到在线系统}
\paragraph{工业部署的验收指标。}
验收不只看离线分数，还应包括从输入到动作可用的总延迟、可用率、告警频率、平均提前量、误停机率、人工复核负担和模型回滚路径。模型上线前应进行影子运行，让模型产生建议但不直接控制设备；通过一段时间验证告警稳定性后，再逐步扩大影响范围。

\paragraph{将验收指标写成可执行约定。}
例如可预先规定：数据到达至建议可用的总延迟不超过2秒，超过该时间则视为超时而不是算作及时正确；通信中断时回退到既有规则；连续多次相同告警合并为一条事件。影子运行中记录这些规则是否被满足，不把模型建议送入控制器。具体门槛应由现场风险和处置能力确定，此处数字只演示怎样把“响应快、告警少”转成可核查条目。

\section{案例分析：剩余寿命预测}
\paragraph{问题定义与数据边界。}
设设备在时刻$t$的观测为$X_{\leq t}$，剩余寿命为$RUL_t$。可以将问题写成点预测$\hat r_t=f_\theta(X_{\leq t})$，也可以预测条件分布$p(r\mid X_{\leq t})$。后者能表达不同设备在相同观测下的差异，但评价要同时关注误差、区间覆盖率和提前量。若模型在寿命末期才准确，可能对维护决策没有价值。

\paragraph{寿命标签与删失似然。}
设从共同起点计算的失效时间为$T$，记录截止时间为$C$，观测到$u=\min(T,C)$与失效指示$\delta=\mathbb 1\{T\leq C\}$。对给定特征$x$，模型定义密度$f_\theta(t\mid x)$和生存函数$S_\theta(t\mid x)=P(T>t\mid x)$。若给定$x$后删失与失效独立，且忽略与$\theta$无关的删失机制，则单台设备的似然贡献是
\begin{equation}
 L_i(\theta)=f_\theta(u_i\mid x_i)^{\delta_i}
 S_\theta(u_i\mid x_i)^{1-\delta_i}.\label{eq:data-censored-likelihood}
\end{equation}
式\eqref{eq:data-censored-likelihood}中，已失效设备贡献准确失效时刻的密度，未失效设备只说明寿命大于$u_i$，因此贡献尾部概率，不能赋予“寿命等于截止时间”的标签。独立设备的负对数似然是$-\sum_i[\delta_i\log f_\theta(u_i\mid x_i)+(1-\delta_i)\log S_\theta(u_i\mid x_i)]$。

在固定协变量模型中，已运行到$t$仍存活的设备，其剩余寿命尾概率为
\begin{equation}
 P(T-t>r\mid T>t,x)=\frac{S_\theta(t+r\mid x)}{S_\theta(t\mid x)},\qquad r\geq0.
\end{equation}
该式要求$S_\theta(t\mid x)>0$。由于$r\geq0$时事件$\{T>t+r\}$包含于$\{T>t\}$，条件概率的分子就是$S_\theta(t+r\mid x)$，由此得到比值。若均值存在，条件期望剩余寿命等于该尾概率对$r$从零到无穷的积分。这个比值明确体现“已活到当前”的信息。若输入是随时间变化的完整历史，则需要相应动态条件模型，不能把未来传感器值塞进固定$x$。预防性维修可能依赖潜在风险，导致信息性删失；此时独立删失假设应重新检查，单用上述似然会偏差。

\paragraph{一台失效、一台尚未失效时怎样贡献证据。}
以指数寿命模型为例，故障率$\lambda>0$的单位为每小时，$f(t)=\lambda e^{-\lambda t}$、$S(t)=e^{-\lambda t}$。设备A在5小时失效，设备B记录到8小时仍存活，二者独立时似然为$(\lambda e^{-5\lambda})e^{-8\lambda}=\lambda e^{-13\lambda}$。对数似然导数为$1/\lambda-13$，最大似然解为$1/13$每小时。若错误地把B也当作8小时失效，就会得到$2/13$，人为提高故障率。这里密度$f(5)$不是“恰好5小时失效”的点概率；连续时间中，一个很小区间内的概率近似为密度乘区间长度。

这个指数模型的$S(t+r)/S(t)=e^{-\lambda r}$与$t$无关，即已经运行多久不改变剩余寿命分布。这是很强的无记忆假设，不适合所有老化设备；例子用于解释删失的贡献，而不是推荐所有寿命都服从指数分布。

\section{数据时序、缺失与来源记录的深入分析}
\paragraph{事件时间与处理时间。}
流式系统应区分事件发生、采集、到达、处理与最终可用的时刻。处理时间是系统实际执行该步计算的时刻，不一定等于入库或到达时间；排队与后续转换还可能推迟特征可查询的时间。标签窗口按事件时间定义，在线输入则检查最终可用时刻$a_j\leq t$。回放时也必须按这些可用时间释放记录，不能一次加载完整历史后只按事件时间截取。

\paragraph{窗口化的统计代价。}
重叠窗口可以增加样本数量，但相邻窗口高度相关，不能把它们当作独立样本计算置信区间。若窗口长度为$L$、步长为$s$，在$s\ll L$时有效样本量远小于窗口数量。设备级切分应在窗口化之前或以设备事件为单位进行，避免同一故障轨迹被拆到多个集合。

\paragraph{有效样本量的计算依据。}
设窗口指标$Z_1,\ldots,Z_N$弱平稳，$\Var(Z_i)=\sigma_Z^2$，滞后$h$相关系数为$\rho_h$。对均值展开协方差和可得
\begin{align}
 \Var(\bar Z)
 &=\frac1{N^2}\left[N\sigma_Z^2+2\sum_{h=1}^{N-1}(N-h)\sigma_Z^2\rho_h\right]\\
 &=\frac{\sigma_Z^2}{N}\left[1+2\sum_{h=1}^{N-1}(1-h/N)\rho_h\right].\label{eq:data-correlated-mean-variance}
\end{align}
把式\eqref{eq:data-correlated-mean-variance}中的括号记为$\tau_N$，按独立均值方差定义$N_{\rm eff}=N/\tau_N$。这是一种针对特定统计量的等效精度，不是数据集固有的唯一属性，也不能只由窗口长度和步长精确确定。正相关常使有效数量下降，负相关则可能相反；非平稳退化序列不能直接套用此估计。实践中以设备或独立运行过程做成组重采样，通常比把相关系数机械外推到无穷更可解释。

\paragraph{相关窗口为何不能充当更多独立设备。}
若两个窗口指标方差都为1、相关系数为$0.8$，它们平均值的方差为$(1+1+2\times0.8)/4=0.9$，而两个独立指标平均的方差为$0.5$。按同一精度折算，两个相关窗口只有$2/(1+0.8)\approx1.11$个独立指标的信息量。这里“弱平稳”要求均值和方差不随时刻改变，协方差只依赖时间间隔；持续退化造成均值变化时，应先检查这一前提，而非机械套公式。

\paragraph{缺失机制与敏感性分析。}
缺失值本身可能携带信息：传感器在高温时掉线，或维护期间停止采集，都不是完全随机缺失。模型可以使用缺失指示变量，但这也可能学到设备或时间的伪相关。应比较删除、简单填补、模型填补和显式缺失建模，并在不同缺失率下报告性能曲线。

\paragraph{异常检测与修复边界。}
统计异常、物理异常和业务异常需要不同处理。一个高振幅冲击可能是故障，也可能是传感器饱和；自动裁剪会消除最有价值的信号。建议同时保留原始值、质量标记和修正值，让模型或分析者能够区分“观测异常”和“处理异常”。

\paragraph{数据版本与来源记录。}
保存一版数据集时，应同时保存它对应的原始文件、转换脚本、特征配置、标签规则和切分清单。数据来源与处理记录指的是这条来路及处理过程。例如看到一个振动均方根字段，应能追查它来自哪台设备、哪段波形，用哪个脚本算出，是否误用了预测之后的记录，以及线上哪个版本读取了它。这样发现错误告警时，才能沿处理过程定位问题，而不只是重新训练一次。

\section{实验归因与部署反馈}
\paragraph{实验矩阵与消融。}
主实验应先固定切分和评价规则，再逐项改变模型或特征。消融实验需要回答机制问题，例如去掉频域特征后是否仍能识别故障、去掉缺失掩码后是否对掉线更脆弱、缩短窗口后提前量如何变化。只罗列大量模型名称而没有机制消融，不能说明改进来自何处。

\paragraph{消融中的交互作用与配对。}
若同时研究频域特征$A$和缺失掩码$B$，设$r_{ab}$是在相同设备、切分和随机种子设置下的风险，$a,b\in\{0,1\}$表示是否启用。$A$在无$B$时的作用为$r_{10}-r_{00}$，在有$B$时为$r_{11}-r_{01}$，两者之差
\begin{equation}
 I_{AB}=r_{11}-r_{01}-r_{10}+r_{00}
\end{equation}
衡量交互作用。因此“一次只改变一个因素”适合局部归因，但不能证明组合收益是两项收益的相加；必要时应使用完整的四种设置。各设置应在同一测试设备上配对比较，并分别报告跨设备与训练种子波动。更多种子只能衡量算法随机性，不能替代更多独立设备。

\paragraph{用四个风险值读交互作用。}
若$r_{00}=10,r_{10}=8,r_{01}=9,r_{11}=6$，则没有掩码时加频域特征使风险下降2，有掩码时下降3，$I_{AB}=6-9-8+10=-1$，表示两者组合比风险加性预测的7再低1。这只是当前配对实验中的交互估计，还需设备级不确定性，不能把负号本身当作机制已经证实。

\paragraph{从窗口告警到事件验收。}
设第$e$次故障发生在$T_e$，预先规定有效告警窗口$[T_e-H,T_e-h]$，其中$H\geq h\geq0$，$h$为最低操作提前量。若窗口内至少一次合法告警，则该事件记为命中；首个符合规则的告警时刻为$A_e$，提前量为$T_e-A_e$。事件召回率是命中事件数除以可完整评价的故障事件数，而非阳性窗口命中比例。还应规定重复告警合并、静默期和告警到事件的匹配规则。

例如两次可完整评价的故障只命中一次，事件召回率为$1/2$；即使该命中事件前有十个阳性窗口告警，也不能计成十次故障命中。若没有可完整评价的故障，事件召回率无定义；若没有命中事件，命中条件下的平均提前量也无定义。报告时应保留相应计数并标注不可计算，而不是填成满分或零提前量。

误报可按单位无故障暴露时间计为$N_{\rm false}/E_{\rm normal}$，从而区分运行一小时与运行一千小时的负担；只有$E_{\rm normal}>0$时才能计算该比率，并应同时报告暴露时间及其单位。提前量均值若只在命中事件上计算，会遗漏漏检事件，故必须与事件召回率一起报告。服务延迟则应从最新合法观测到动作可用的整个过程测量，不能只报神经网络一次前向传播的时间。

\paragraph{自测：没有失效也有信息吗？}
沿用指数寿命模型，两台独立设备分别观察到4小时和6小时，均未失效。写出似然；能否像出现一次故障时一样求得正的内部最大似然估计？

\textit{参考解。}似然为$e^{-4\lambda}e^{-6\lambda}=e^{-10\lambda}$。它随$\lambda>0$严格下降，只有$\lambda\to0$时逼近上确界，没有正的内部最大值。记录提供了存活证据，却没有精确识别一个正故障率；不能把截止时刻当作失效时刻来制造两个故障。

\paragraph{从离线到影子运行。}
上线前先回放历史数据，再进行不影响控制回路的影子运行。影子阶段检查数据延迟、字段缺失、告警频率、模型版本和人工反馈。通过验收后才逐步开放建议或控制权限，并保留回滚阈值。上线监控应同时关注输入漂移、输出漂移、延迟到达的真实标签和服务资源。标签到达后才能回查漏报、误报与校准；这不同于模型计算的后验概率。

本章的预处理、窗口和数据质量规则可用下列官方文档交叉检查，尤其注意变换的拟合时点。
\section*{辅助学习材料}
\begingroup
\emergencystretch=3em
\begin{itemize}
\item \textbf{Preprocessing data}；作者/平台：scikit-learn Developers；英文；官方文档。阅读 标准化、缩放与类别编码小节，帮助理解本章的缺失机制、量纲和训练侧统计量。\href{https://scikit-learn.org/stable/modules/preprocessing.html}{在线阅读}。
\item \textbf{Signal Processing}；作者/平台：SciPy Developers；英文；官方教程。重点阅读 filtering、Fourier transform 和 resampling 相关导航，帮助把本章的采样、混叠、频域特征与时间窗口连接起来。\href{https://docs.scipy.org/doc/scipy/tutorial/signal.html}{在线阅读}。
\item \textbf{【机器学习】Cross-Validation（交叉验证）详解}；知乎；中文解说。用于本章实验设计与数据划分，重点检查分组和时间边界；不替代缺失机制与信号处理内容。标题与主题据必应索引，原页受限。\href{https://zhuanlan.zhihu.com/p/24825503}{在线阅读}。
\item \textbf{交叉验证（Cross-Validation）}；CSDN；中文博客。用于划分流程检查；预处理必须仅在各训练折拟合，工业数据还需独立处理设备和时间相关性。标题与主题据必应索引，原页受限。\href{https://blog.csdn.net/weixin\_42691585/article/details/113971857}{在线阅读}。
\item \textbf{13-数据预处理-标准化}；黑马程序员；bilibili；中文课程分集。选看《黑马程序员3天快速入门python机器学习》第13集，结合第12集归一化比较尺度变换；对照本章要求，仅在训练集拟合变换参数。2026年10月4日公开元数据所列整套课程累计播放1,793,918次（不是本分集计数）；已核对目录与计数，未逐帧观看。\href{https://www.bilibili.com/video/BV1nt411r7tj/?p=13}{观看分集}。
\end{itemize}
\endgroup
\paragraph{本章小结。}
工业数据工程把现实过程变成时间边界清楚、版本可追溯、统计依赖可解释的学习样本。模型性能的上限往往由数据管道决定；越复杂的模型，越需要可靠的数据来源与处理记录和部署验收来支撑其结论。

时间集合规定能用什么信息，缺失和删失规定我们究竟观察到了什么，成组评价与配对消融规定改进可以归因到哪里。它们共同完成第一部分的学习任务设置：不再只交付一个拟合函数，而是交付一条从记录到预测再到维护动作的可复核的处理过程。

记录整理正确以后，设备数量有限仍会使故障率和回归系数难以精确估计，当前隐状态也可能有多种解释。下一章用先验、似然和后验描述这些不确定性。输入若混入未来信息，概率推断也无法补救；输入使用正确时，它才有助于判断哪些结论较稳定，哪些还需补充数据或请人复核。
